\chapter{Mathematical model and modelling assumptions}
\label{ch:model}

\section{Geometry and reference configuration}
\label{sec:domains}

The configuration is two-dimensional. Reference coordinates $X=(X,Y)$ label material points of the
solid and the undeformed fluid domain. The fluid reference domain,
$\Om_f^0=\{(X,Y):-L_e\le X\le L+L_o,\;0\le Y\le H_0\}$,
consists of a rigid entrance section $\Om_e=(X\in[-L_e,0])$, a section $\Om_d=(X\in[0,L])$ facing
the deformable wall, and a rigid exit section $\Om_o=(X\in[L,L+L_o])$, all of undeformed height
$H_0$. The only deformable body is the solid strip
$\Om_s^0=[0,L]\times[H_0,H_0+d]$
of thickness $d$ spanning the full out-of-plane depth. Its lower face $Y=H_0$ shares the fluid
boundary; its upper face $Y=H_0+d$ receives the imposed thermal load. The strip is embedded in a
rigid casing and is pinned at both ends, $X=0$ and $X=L$: the displacement vanishes there while the
bending moment vanishes (Assumption~A9). Figure~\ref{fig:geometry} shows the configuration,
Table~\ref{tab:geometry} records the design point.

\begin{figure}[t]
\centering
% fig_geometry.tex -- body content of Figure 1 (scaled layout, not to absolute scale)
\begin{tikzpicture}[x=1pt,y=1pt,>=Stealth,font=\small]
% ---------------- geometry (points as described in caption) ----------------
\def\Le{36}\def\L{180}\def\Lo{36}          % entry, heated, exit (horizontal pts)
\def\Hf{22}                                 % channel height (exaggerated for readability)
\def\ds{16}                                 % wall thickness. origin of fluid at lower-left of entry zone
\coordinate (O) at (0,0);
% fluid channel
\fill[blue!10] (O) rectangle ++(\Le+\L+\Lo,\Hf);
\draw[thick] (0,0) -- ++(\Le+\L+\Lo,0);
\draw[thick] (0,\Hf) -- ++(\Le,0) coordinate (L1);
\draw[thick] (\Le+\L,\Hf) -- ++(\Lo,0);
% flexible wall strip (heated segment), deformed bow shown by parabola through points
\coordinate (A) at (\Le,\Hf);
\coordinate (B) at (\Le+\L,\Hf);
\coordinate (At) at (\Le,\Hf+\ds);
\coordinate (Bt) at (\Le+\L,\Hf+\ds);
\fill[orange!25] (At) -- (Bt) -- (B) .. controls (\Le+\L/2,\Hf-12) .. (A) -- cycle;
\draw[thick] (A) .. controls (\Le+\L/2,\Hf-12) .. (B);
\draw[thick] (At) -- (Bt);
\draw[thick] (A) -- (At); \draw[thick] (B) -- (Bt);
% heat flux arrows on top of heated wall
\foreach \x in {0.12,0.30,0.5,0.7,0.88}{
  \draw[->,red!70!black,thick] (\Le+\x*\L,\Hf+\ds+14) -- (\Le+\x*\L,\Hf+\ds+2);}
\node[red!70!black] at (\Le+\L/2,\Hf+\ds+20) {$q^{''}$ prescribed, $\Ga_h$};
% adiabatic labels
\node[align=center] at (\Le/2,\Hf+9) {adiabatic rigid lid\\(solid extension)};
\node[align=center] at (\Le+\L+\Lo/2,\Hf+9) {adiabatic};
% inlet/outlet arrows
\draw[->,thick] (-16,\Hf/2) -- (0,\Hf/2) node[midway,above left=1pt] {$\bu=(U,0)$,\; $T_\mathrm{in}$};
\draw[<-,thick] (\Le+\L+\Lo,\Hf/2) -- ++(16,0);
\node[below] at (\Le+\L+\Lo+8,\Hf/2-4) {outlet, $p=p_\mathrm{out}$};
% interface and wall-motion annotation
\draw[<->,blue!60!black] (\Le+\L/2,\Hf-12) -- (\Le+\L/2,\Hf-1);
\node[blue!60!black,anchor=west] at (\Le+\L/2+4,\Hf-6) {$w(x)$, $\Ga$};
% dimension lines
\draw[<->] (\Le,-8) -- ++(\L,0) node[midway,below] {$L$, heated flexible wall};
\draw[<->] (0,-8) -- ++(\Le,0) node[midway,below] {entrance $L_e$};
\draw[<->] (\Le+\L,-8) -- ++(\Lo,0) node[midway,below] {exit $L_o$};
\draw[<->] (-10,0) -- ++(0,\Hf) node[midway,left] {$H_0$};
\draw[<->] (\Le+\L+6,\Hf) -- ++(0,\ds) node[midway,right] {$d$};
% pinned supports
\foreach \p in {(A),(B)}{\draw \p -- ++(-4,-6) -- ++(8,0) -- cycle;}
\node[anchor=west] at (\Le+\L+4,-6) {pinned ends (parameter-locked, register D-4)};
\end{tikzpicture}

\caption{Configuration of the coupled problem (schematic, not to scale). Rigid entrance and exit
sections of the channel; a thermoelastic segment of length $L$ and thickness $d$, pinned at its
ends and heated on its outer face $\Ga_h^0$; the conjugate interface $\Ga_{\mathrm{if}}$ between
fluid and solid; a rigid, adiabatic lower wall. The interface deflection $w$ is measured positive
away from the channel (Figure~\ref{fig:signs}).}
\label{fig:geometry}
\end{figure}

\begin{table}[t]
\centering\small
\caption{Geometry and operating point used as the baseline design. Values marked \emph{design} are
free choices of the study; the derived hydraulic diameter uses the slit value $D_h=2H_0$.}
\label{tab:geometry}
\begin{tabular}{p{4.0cm}p{1.7cm}p{4.2cm}p{2.6cm}}
\toprule
Quantity & Symbol & Value & Basis \\
\midrule
Channel height & $H_0$ & \SI{5.0e-4}{\metre} & design \\
Hydraulic diameter & $D_h=2H_0$ & \SI{1.0e-3}{\metre} & derived \\
Heated, deformable length & $L$ & \SI{3.0e-2}{\metre} & design \\
Entrance / exit lengths & $L_e,L_o$ & \SI{1.0e-2}{\metre}; \SI{2.0e-2}{\metre} & design \\
Wall thickness & $d$ & \SI{3.0e-3}{\metre} & design \\
Mean inlet velocity & $U$ & \SI{5.0e-2}{\metre\per\second} & design \\
Imposed heat flux (baseline regime) & $q^{''}$ & \SI{5.0e3}{\watt\per\square\metre} & design \\
Inlet and reference temperatures & $T_{\mathrm{in}},T_{\mathrm{ref}}$ & to be fixed at data freeze & temperature differences only enter \\
\bottomrule
\end{tabular}
\end{table}

The deformed interface is defined by the solid displacement field $\bd$ on the reference interface,
\begin{equation}
\Ga_{\mathrm{if}}=\{x=X+\bd(X):\;X\in[0,L]\times\{H_0\}\},
\end{equation}
and the steady fluid domain $\Om_f$ is obtained by extending the interface displacement into
$\Om_f^0$ through a mesh-extension map $\widehat{\mathbf w}$ specified in
Section~\ref{sec:ale}. The fluid equations are posed on $\Om_f$; the solid equations on the
reference configuration, consistent with small strains (Assumption~A4).

\section{Modelling assumptions}
\label{sec:assumptions}
Each assumption is stated with its justification and its range of validity; the quantitative checks
supporting the numerical bounds are carried out in Chapter~\ref{ch:baseline}.

\begin{itemize}\itemsep2pt
\item[\textbf{A1}] \emph{Two-dimensional geometry and plane strain.} The out-of-plane dimension is
large compared with $H_0$ and $d$, and loading is uniform along it, so the flow and temperature
fields are two-dimensional and the solid is in plane strain, $\eps_{zz}=0$. Plane strain, rather
than plane stress, is the closure appropriate to a wide bonded wall
\cite{boley1960theory,timoshenko1970elasticity}; its consequences for the constitutive reduction
are derived in Chapter~\ref{ch:derivations}.
\item[\textbf{A2}] \emph{Newtonian, constant-property fluid.} The working fluid is water at near-ambient
temperature, with constant density $\rho$, viscosity $\mu$, specific heat $c_p$, and conductivity
$k_f$. The self-consistency of constant properties is established a posteriori: the computed
fluid-side temperature excursions at baseline are of order $1$--$3$ K
(Section~\ref{sec:baselinecalc}), over which the property variation remains below a few per cent.
The influence of buoyancy is examined separately through the Richardson number
(Sections~\ref{sec:buoyancy-baseline}, \ref{sec:limitations}).
\item[\textbf{A3}] \emph{Homogeneous isotropic linear thermoelastic solid.} The wall is a PEEK-class
polymer with constant $E$, $\nu$, $\alpha$, $k_s$, $\rho_s$, $c_s$. Deviations (grade-to-grade
variability, mild anisotropy of extruded stock) are treated by the sensitivity requirements of
Chapter~\ref{ch:validation}; the numerical values used below are datasheet-representative and are
to be frozen against a specific datasheet before implementation (Appendix~B, item C-6).
\item[\textbf{A4}] \emph{Small strains, moderate rotations.} The strain measures must remain small:
the bending fibre strain, the arc-stretching membrane strain, and the total mechanical surface
strain. Linear kinematics alone is \emph{not} sufficient for the pinned, axially constrained strip,
because the thermal membrane force interacts with the transverse deflection at leading order; the
adopted kinematics is therefore the reduced (frozen-force) von K\'arm\'an beam--column level of
Section~\ref{sec:deriv:beamcolumn}, which reduces to the linear relation when the membrane force
is negligible. At the baseline the bending fibre strain is $9.9\times10^{-4}$, the stress-producing
surface strain is at most $3.25\times10^{-3}$ and the kinematic envelope $e_0+\kappa_{\mathrm{geo}}
d/2$ is $5.35\times10^{-3}$, the arc-stretching term is
$\approx1.39\times10^{-4}$, and the maximum slope is $0.027$~rad
(Section~\ref{sec:baselinecalc}); all are far below the order-one failure limits of the
hypothesis. Geometric admissibility is expressed through these strain and slope measures together
with a mesh bound $|w|/H_0\le0.5$ with positive element Jacobians, rather than through a purely
geometric deflection ratio alone (Appendix~B, register D-11, D-15).
\item[\textbf{A5}] \emph{Steady flow and energy; quasi-static solid mechanics.} The structural wave
time across the wall is seven orders of magnitude shorter than the slowest transport time
(Section~\ref{sec:steadytimes}), so the mechanical problem is quasi-static and inertia is omitted
from \eqref{eq:equilibrium}.
\item[\textbf{A6}] \emph{Laminar flow.} The design keeps $\mathrm{Re}_{D_h}\le100$ (baseline value $50$),
well below the transition range for plane channels, which justifies \eqref{eq:momentum}.
\item[\textbf{A7}] \emph{Negligible viscous dissipation.} The Brinkman number at baseline is of order
$10^{-6}$ (Chapter~\ref{ch:baseline}); the dissipation term is retained as a re-inclusion option
with the rule stated at \eqref{eq:fluidenergy}.
\item[\textbf{A8}] \emph{Perfect thermal contact at the interface.} No contact resistance is included at
$\Ga_{\mathrm{if}}$; temperature and heat flux are continuous
(\eqref{eq:iftemp}--\eqref{eq:ifflux}).
\item[\textbf{A9}] \emph{Pinned supports and rigid casing.} Each strip end is a \emph{pin} in the
structural sense: the end section may rotate about the transverse axis, carries no applied moment,
and does not translate axially on average, so the mean axial expansion of the strip is blocked.
On each end edge ($X=0,L$) the conditions are therefore $d_2=0$, $\langle d_1\rangle=0$ taken as the
mean over the edge, zero applied moment (a natural condition), and zero heat flux. They are
\emph{not} $d_1=d_2=0$ pointwise: prescribing both displacement components on the whole edge is a
\emph{clamp}, which forbids end rotation, contradicts the moment-free condition and suppresses the
bowing altogether, since a clamped strip at uniform $\kappa_{\mathrm{th}}$ has $w\equiv0$
(Section~\ref{sec:deriv:curvature}). An independent plane-strain Q8 computation, reported in
Appendix~B (register D-26), gives $1.0000$ of the beam value $\kappa_{\mathrm{th}}L^2/8$ for point
pins, $0.998$ for a full-edge vertical support with a free axial mean, and $-0.02$ to $-0.06$ for
$d_1=d_2=0$ on the whole edge. That comparison is a finite-mesh result on one $24$-element mesh,
not a converged strip solution (register D-26); it shows that the full-edge vertical support costs
$0.2\%$, while the axial restraint of a clamp destroys the mechanism.
The surrounding casing is taken rigid and adiabatic, and its interior pressure equals the outlet
datum of condition~B2, which fixes the load reference of Section~\ref{sec:deriv:pressure}.
\item[\textbf{A10}] \emph{Uniform imposed flux as the baseline thermal load (regime N).} A uniform
$q^{''}$ models a film heater and fixes the total heat input per unit depth, $Q'=q^{''}L$, which is
the quantity entering the global enthalpy balance (Section~\ref{sec:quality}). A second regime
(regime D, imposed outer-face temperature $T_h$) is retained as a separately labelled study;
comparisons across regimes are meaningful only at matched total heat input.
\end{itemize}

\section{Sign and traction conventions}
\label{sec:signs}
A single interface normal $\bn$ is used throughout, directed from the fluid into the solid.
Pressure on the wall is positive in compression; Cauchy stress in the solid is positive in tension.
The interface deflection is denoted $w$ and is measured \emph{from the reference interface away from
the channel} (hence $w=+d_2$ in terms of the upward-positive solid displacement $d_2$). Wall curvature
is taken positive when the wall bows away from the channel; within the small-slope assumption this gives
the design relation $\kappa=-w^{''}$, and a constant eigen-curvature $\kappa_{\mathrm{th}}>0$
produces $w=\kappa_{\mathrm{th}}x(L-x)/2\ge0$ (Section~\ref{sec:deriv:deflection}).
The pairing is not a matter of taste. A pinned end requires the physical first moment
$\int\sigma_{xx}y\,\mathrm dy$ to vanish there; with the opposite curvature pairing that residual is
$-27$~N\,m at $x=0$ --- some four times the $6.4$~N\,m the strip actually carries at mid-span --- so
only the pairing stated above is compatible with the pinned end (Appendix~B, register D-21).
Fluid heat flux along $+\bn$ carries heat out of the fluid, so the flux \emph{into} the fluid is
$+k_f\nabla T_f\cdot\bn$ and is positive on the heated wall. Using one normal throughout makes the
traction equality \eqref{eq:iftraction} and the flux equality \eqref{eq:ifflux} single-valued and
gives the two subdomains opposite boundary-term signs. That convention is bookkeeping: it does not
by itself make any interface integral cancel, and Section~\ref{sec:coupling} states what does and
does not cancel and under which conditions.

\begin{figure}[t]
\centering
% fig_signs.tex -- body content of Figure 2 (sign/traction conventions at the interface)
\begin{tikzpicture}[x=1pt,y=1pt,>=Stealth,font=\small]
% fluid (below) / solid (above) separated by interface
\fill[blue!10] (-10,-44) rectangle (190,0);
\fill[orange!25] (-10,0) rectangle (190,26);
\node at (40,-14) {fluid $\Om_f$};
\node at (150,17) {solid $\Om_s^{0}$};
% interface
\draw[very thick] (-10,0) -- (190,0) node[midway,above=2pt] {$\Ga$};
% normal convention: n points out of the fluid, i.e. up and into the solid
\draw[->,thick,blue!60!black] (55,0) -- ++(0,22) node[midway,right=2pt] {$\bn$ (out of fluid)};
% by (5.1) with the single normal n (out of fluid, into solid): the force ON the solid is
% -sigma_f n (pressure part +p n, pointing up, away from the fluid); sigma_f n is the reaction on the fluid
\draw[->,thick,red!70!black] (95,0) -- ++(0,20) node[midway,left] {$-\bsig\bn$: force ON solid};
\draw[->,thick,red!70!black] (95,0) -- ++(0,-20) node[midway,right] {$\bsig\bn$: reaction traction on fluid};
% wall shear sense
\draw[->,thick] (130,-8) -- ++(26,0) node[midway,below] {$\tau_w$ (flow direction)};
% flux continuity
\node[align=left] at (105,-32) {flux: $-k_f\partial_nT_f=-k_s\partial_nT_s$ on $\Ga$ (single normal convention)};
\end{tikzpicture}

\caption{Interface conventions. A single unit normal $\bn$ (fluid $\to$ solid) is used in the
traction equality \eqref{eq:iftraction}, in the flux equality \eqref{eq:ifflux}, and in all
variational boundary terms of Section~\ref{sec:coupling}.}
\label{fig:signs}
\end{figure}

\section{Material data}
\label{sec:props}
Table~\ref{tab:props} lists the property values used in all calculations. Fluid values are standard
near-ambient water data. Solid values are representative of unfilled PEEK; they are marked
accordingly, because they must be frozen against a chosen datasheet before the numerical
implementation (Appendix~B, item C-6), and the derived baseline quantities that depend on them are
reported with that dependency stated.

\begin{longtable}{lllll}
\caption{Material properties used in the calculations. ``freeze pending'' = datasheet value to be
fixed before implementation (Appendix~B, item C-6).}\label{tab:props}\\
\toprule
Property & Symbol & Value & Units & Remark \\
\midrule
\endfirsthead
\toprule
Property & Symbol & Value & Units & Remark \\
\midrule
\endhead
Water density & $\rho$ & $1.0\times10^{3}$ & \si{\kilogram\per\cubic\metre} & standard value \\
Water dynamic viscosity & $\mu$ & $1.0\times10^{-3}$ & \si{\pascal\second} & standard value \\
Water specific heat & $c_p$ & $4.18\times10^{3}$ & \si{\joule\per\kilogram\per\kelvin} & standard value \\
Water conductivity & $k_f$ & $0.60$ & \si{\watt\per\metre\per\kelvin} & standard value \\
Water expansion coefficient & $\beta_f$ & $2.1\times10^{-4}$ & \si{\per\kelvin} & freeze pending \\
Water kinematic viscosity & $\nu_{\mathrm{kin}}=\mu/\rho$ & $1.0\times10^{-6}$ & \si{\square\metre\per\second} & derived \\
PEEK Young's modulus & $E$ & $3.6\times10^{9}$ & \si{\pascal} & freeze pending \\
PEEK Poisson ratio & $\nu$ & $0.40$ & -- & freeze pending \\
PEEK expansion coefficient & $\alpha$ & $5.0\times10^{-5}$ & \si{\per\kelvin} & freeze pending \\
PEEK conductivity & $k_s$ & $0.25$ & \si{\watt\per\metre\per\kelvin} & freeze pending \\
PEEK density & $\rho_s$ & $1.3\times10^{3}$ & \si{\kilogram\per\cubic\metre} & freeze pending \\
PEEK specific heat & $c_s$ & $1.1\times10^{3}$ & \si{\joule\per\kilogram\per\kelvin} & freeze pending \\
\bottomrule
\end{longtable}

\section{Notation}
Current and reference coordinates are $x=(x,y)$ and $X=(X,Y)$. Fluid unknowns are the velocity
$\bu=(u,v)$, pressure $p$, and temperature $T_f$; solid unknowns are displacement
$\bd=(d_1,d_2)$, temperature $T_s$, stress $\bsig$, and strain $\beps$;
$\theta=T_s-T_{\mathrm{ref}}$ denotes the temperature rise in the solid. The Lam\'e parameters are
written $\lambda$ and $\mu_L$, where $\mu_L=G$ is the shear modulus (the subscript distinguishes it
from the fluid viscosity). The bulk modulus is $K=\lambda+\tfrac23\mu_L=E/[3(1-2\nu)]$, and the
thermoelastic pressure coefficient is
$\beta=(3\lambda+2\mu_L)\alpha=3K\alpha=E\alpha/(1-2\nu)$
\cite{boley1960theory}. A complete symbol table with units is given in Appendix~\ref{app:symbols}.
